\chapter{ФОРМУЛЮВАННЯ ОДНОРІДНИХ ЗАДАЧ У М'ЯКИХ ТКАНИНАХ ТІЛА і МЕТОДІВ ЇХ РОЗВ'ЯЗУВАННЯ}

\section{Стаціонарна задача теплопровідності}

Серцево-судинна система людини є складною. У зв'язку із цим є два основні підходи до моделювання цієї системи в 
організмі людини \cite{kutz-zhu-heat-transfer-biological-systems}:
\begin{itemize}
    \item Моделі суцільного середовища, де береться усереднений ефект від потоку крові через об'єм у модельованих м'яких
    тканинах. 
    \item Судинні моделі, у яких моделюються головні судини, які впливають на температурні процеси органу чи органів, 
    які досліджуються.
\end{itemize}
\noindent Удосконаленням медичних інструментів полегшує створення судинних моделей, які є складнішими, але дають точніші
результати. 

У цій праці будемо використовувати модель суцільного середовища. Вона базується на статті Пеннеса
\cite{thermal-process-tissue-pennes}, яка наводить вже класичне рівняння біотеплового процесу у людському організмі. Ця
стаття, була опублікована у 1948 році, але до тепер залишається актуальною для більшості досліджень. Зокрема це рівняння
використовується для оцінки ракових утворень у легенях \cite{lung-tumor-thermal-conductivity}.

Нехай нам задана область $\Omega$. Границю області позначимо $\Gamma$, $\Gamma = \partial \Omega$. Стаціонарне 
температурне поле є розв'язком такого рівняння теплопровідності:

\begin{equation}
    \label{eqn:theory_thermo_pennes_eqn}
    \nabla \cdot \lambda_\theta \nabla \theta(x) + \omega \rho_b C_b (\theta_c - \theta(x)) + q_m  = 0, x \in \Omega
\end{equation}

\noindent де $\lambda_\theta$ -- коефіцієнт теплопровідності, $\theta$ -- шукана температура в області $\Omega$, 
$\rho_b$ -- густина крові, $\omega$ -- перфузія крові через орган, $C_b$ -- питома теплоємність крові, 
$\theta_c$ -- внутрішня температура тіла, $q_m$ -- метаболічне тепло у тканині. На межі задамо наступні умови:

\begin{equation}
    \label{eqn:theory_thermo_cond_dirichlet}
    \theta(x) = \theta_{\Gamma_D}(x), x \in \Gamma_D 
\end{equation}

\begin{equation}
    \label{eqn:theory_thermo_cond_neumann}
    q(x) = q_{\Gamma_N}(x), x \in \Gamma_N
\end{equation}

\begin{equation}
    \label{eqn:theory_thermo_cond_robin}
    -\lambda_\theta q(x) = h(\theta(x) - \theta_e), x \in \Gamma_R,
\end{equation}

\noindent де $q(x) = \frac{\partial \theta(x)}{\partial n}$ -- тепловий потік, $n_i$ -- компоненти вектора зовнішньої 
нормалі до межі $\Gamma$, $\theta_e$ -- температура середовища, $h$ -- коефіцієнт тепловіддачі. Граничні умови 
\ref{eqn:theory_thermo_cond_dirichlet}, \ref{eqn:theory_thermo_cond_neumann} і \ref{eqn:theory_thermo_cond_robin} -- це 
умови Діріхле, Неймана і Робіна відповідно. Умова Робіна описує тепловий обмін тіла із навколишнім середовищем.

\section{Стаціонарна задача пружності}

Розглянемо область $\Omega$. $u_i(x)$ -- позначимо компоненти вектора переміщень. Компоненти тензора деформацій 
$\varepsilon_{ij}$ визначаються співвідношеннями Коші:  

\begin{equation}
    \label{eqn:theory_elasticity_gauchy_tensor}
    \varepsilon_{ij}(x) = \frac{1}{2}(u_{i,j}(x) + u_{j,i}(x)), x \in \Omega
\end{equation}

Компоненти тензора напружень $\sigma_{ij}$ можна виразити через $\varepsilon_{ij}$ використавши закон Гука:

\begin{equation}
    \label{eqn:theory_elasticity_hooke_law}
    \sigma_{ij}(x) = \lambda \delta_{ij}\varepsilon_{kk}(x) + 2\mu\varepsilon_{ij}(x), x \in \Omega,
\end{equation}

\noindent де $\lambda$ і $\mu$ -- це коефіцієнти Ляме, $\delta_{ij}$ -- символ Кронекера. 

Статична рівновага напружень і моментів виражається таким рівнянням:

\begin{equation}
    \label{eqn:theory_elasticity_sigma_equilibrium}
    \sigma_{ij,i}(x) + b_j(x) = 0, x \in \Omega,
\end{equation}

\noindent де $b_j(x)$ -- компоненти вектора об'ємних напружень. Якщо у рівняння рівноваги 
\ref{eqn:theory_elasticity_sigma_equilibrium} підставити закон Гука \ref{eqn:theory_elasticity_hooke_law} і 
співвідношення Коші \ref{eqn:theory_elasticity_gauchy_tensor}, то одержимо систему диференційних рівнянь у переміщеннях:

\begin{equation}
    \label{eqn:theory_elasticity_displacement_eqn}
    \mu u_{i,jj}(x) + (\lambda + \mu)u_{j,ji} + b_i(x) = 0, x \in \Omega
\end{equation}

\noindent На межі задамо наступні умови:
\begin{equation}
    \label{eqn:theory_elasticity_dirichlet}
    u_i(x) = u_{\Gamma_Di}(x), x \in \Gamma_D
\end{equation}

\begin{equation}
    \label{eqn:theory_elasticity_neumann}
    t_i(x) = t_{\Gamma_Ni}(x), x \in \Gamma_N,
\end{equation}

\noindent де $t_i = \sigma_{ij}n_j$, $n_i$ -- компоненти вектора зовнішньої нормалі до межі $\Gamma$. Граничні умови 
\ref{eqn:theory_elasticity_dirichlet} і \ref{eqn:theory_elasticity_neumann} -- це умови Діріхле і Неймана відповідно.

У практичних задачах замість коефіцієнтів Ляме можуть бути заданими  модуль Юнга і коефіцієнт Пуассона. Для перетворення
цих величин використовуються наступні формули: 

\begin{equation}
    \label{eqn:theory_elasticity_lame_lambda_convertion}
    \lambda = \frac{E\nu}{(1+\nu)(1-2\nu)}
\end{equation}

\begin{equation}
    \label{eqn:theory_elasticity_lame_mu_convertion}
    \mu = \frac{E}{2(1+\nu)}
\end{equation}

\noindent де $E$ -- модуль Юнга, а $\nu$ -- коефіцієнт Пуассона.  

\section{Стаціонарна задача термопружності}

Розглянемо вплив температури на пружність м'яких тканин тіла. Представимо тіло областю $\Omega$. До закону Гука 
\ref{eqn:theory_elasticity_hooke_law} додамо співвідношення Дюгамеля-Неймана і отримаємо:

\begin{equation}
    \label{eqn:theory_thermoelasticity_duhamel_neumann_extension}
    \sigma_{ij}(x) = \lambda \delta_{ij}\varepsilon_{kk}(x) + 2\mu\varepsilon_{ij}(x) - (2\mu + 3\lambda)\alpha\theta, 
        x \in \Omega,
\end{equation}

\noindent де $\alpha$ -- коефіцієнт теплового розширення матеріалу. Температурне поле $\theta$ знайдемо із рівняння 
\ref{eqn:theory_thermo_pennes_eqn}. Запишемо рівняння рівноваги у переміщеннях із використанням 
\ref{eqn:theory_thermoelasticity_duhamel_neumann_extension}:

\begin{equation}
    \label{eqn:theory_thermoelasticity_displacement_eqn}
    \mu u_{i,jj}(x) + (\lambda + \mu)u_{j,ji} - (2\mu + 3\lambda)\alpha\theta + b_i(x) = 0, x \in \Omega
\end{equation}

На границі $\Gamma$ можна задати умови для рівняння Пеннеса \ref{eqn:theory_thermo_cond_dirichlet}, 
\ref{eqn:theory_thermo_cond_neumann} і \ref{eqn:theory_thermo_cond_robin}, а для переміщень і поверхневих напружень 
\ref{eqn:theory_elasticity_dirichlet} і \ref{eqn:theory_elasticity_neumann} відповідно.

\section{Метод граничних елементів}

\section{Метод приграничних елементів}
